\section{Discussion}
\label{sec:discussion}

\paragraph{The role of latent spaces in generalized inversion.}
The flow latent space provides substantial benefit for goal-conditioned inversion ($0.642$ vs.\ $0.523$ in the original space, a $23\%$ improvement) when all methods use the same surrogate. The $4$D manifold reduces the action space and constrains the policy to parameter configurations that the latent model has learned to be physically meaningful, paralleling findings in humanoid motor control~\cite{luo2024pulse} where latent spaces aid generalization across diverse objectives. The flow's bijective mapping also appears to provide implicit regularization: its gap ($\Delta{=}0.012$) is the smallest of any method, suggesting the policy stays in regions where the surrogate is reliable. The VAE latent space is less effective for direct GC-PPO ($0.485$) but produces the best warm-start initialization (WS VAE $0.828$ vs.\ WS Flow $0.773$). One possible explanation is that the VAE's stochastic encoding generates more diverse initial points that seed CMA-ES in different basins, while the flow's deterministic mapping concentrates initializations more narrowly.

\paragraph{Evaluation design and surrogate exploitation.}
Our ablation (Table~\ref{tab:ablation}) shows that using each latent model's own prediction head, rather than the shared MLP surrogate, degrades GC-PPO Flow by $26\%$. This is a practical finding for the broader latent space optimization community: when comparing methods across different representation spaces, the evaluator must be controlled, or performance differences may reflect evaluator accuracy rather than representation quality. More broadly, the surrogate-to-simulator gap $\Delta$ increases monotonically with optimization intensity across all experiments (GC-PPO $\Delta < 0.05$; CMA-ES $540$ $\Delta{=}0.079$; CMA-ES $5{,}000$ $\Delta{=}0.141$). That CMA-ES with $5{,}000$ evaluations actually achieves \emph{lower} recovery than CMA-ES with $540$ evaluations ($0.781$ vs.\ $0.788$) underscores that more search is not always better when optimizing against an imperfect surrogate.

\paragraph{Toward food material identification.}
This work provides the inverse estimation component of a broader pipeline for food material identification. In a complete system, visual tracking methods~\cite{wu2025fruitninja, xie2024physgaussian} would extract behavioral features from video of food manipulation, and the goal-conditioned policy would instantly estimate the corresponding material parameters. The policy's speed (${\sim}10$\,ms per estimate) makes it suitable for real-time or interactive applications. The ${\sim}9$-minute training cost is a one-time investment that generalizes across all future targets.

\paragraph{Limitations.}
We evaluate on a single test case (orange peeling with AnisoMPM); generalization to other food materials and fracture modes is untested. The validation uses $10$ held-out targets with high per-target variance (std $0.084$--$0.162$), limiting statistical power for fine-grained comparisons between methods with similar means. The $9$-dimensional parameter space is modest compared to spatially varying material fields encountered in practice. The warm-start advantage over equal-budget CMA-ES is inconsistent across representations (WS VAE outperforms but WS Flow underperforms CMA-ES $540$), and should not be treated as a universal improvement.

\paragraph{Future work.}
Replacing hand-designed behavioral features with learned representations from video would close the loop between visual observation and inverse estimation. Extending to higher-dimensional parameter spaces would test whether the latent space advantage grows with dimensionality. Incorporating predictive uncertainty into the surrogate could provide a mechanism for detecting and mitigating surrogate exploitation.
